%% 04-reproduction.tex — Step-by-step reproduction

\section{Reproducing the Results}
\label{sec:reproduction}

The following six steps reproduce every result in the portfolio from scratch,
starting from an empty machine with internet access.
Expected total time on modern hardware: approximately two to four hours
for all three census years, or under thirty minutes for 2020 only.

\subsection*{Step 1 --- Clone the Repository}

\begin{verbatim}
git clone https://github.com/giodl73-repo/REDIST
cd REDIST
git checkout v0.2.0      # Pin to the release used in the portfolio
\end{verbatim}

\subsection*{Step 2 --- Run \texttt{bootstrap.sh}}

\begin{verbatim}
./bootstrap.sh            # Linux / macOS
bootstrap.bat             # Windows
\end{verbatim}

\noindent This installs Rust, builds the \texttt{redist} binary, installs
Python dependencies, and runs \texttt{redist doctor} to confirm the
environment is ready.  Expected time: 5--10 minutes.

\subsection*{Step 3 --- Fetch Census Data}

\begin{verbatim}
redist fetch --year 2020      # ~40 GB, allow 30-90 min depending on connection
redist fetch --year 2010
redist fetch --year 2000
\end{verbatim}

\noindent All files are SHA-256 validated on download.
Interrupted fetches resume automatically.

\subsection*{Step 4 --- Build Official Plans}

\begin{verbatim}
redist run --year 2020 --version official_2020    # All 50 states, ~18 min
redist run --year 2010 --version official_2010
redist run --year 2000 --version official_2000
\end{verbatim}

\noindent Output districts are written to
\texttt{outputs/official\_\{year\}/\{year\}/\{state\}/}.
Each district file contains tract-level assignments, population totals,
and the bisection round at which each district was finalized.

\subsection*{Step 5 --- Verify Labels}

\begin{verbatim}
redist label-verify official_2020    # Checks SHA chain for 2020 run
redist label-verify official_2010
redist label-verify official_2000
\end{verbatim}

\noindent This command recomputes the SHA-256 of every config, input file,
and output file, and confirms the chain matches the published values.
If any file was modified---even by a single byte---the command reports
which step of the chain failed.

\subsection*{Step 6 --- Compare SHA Chains}

The published SHA chain values for the official 2020 run are:

\begin{table}[h]
\centering
\caption{Expected SHA-256 chain values for the official 2020 run (placeholders).}
\label{tab:sha}
\begin{tabular}{@{}ll@{}}
\toprule
Chain step & SHA-256 (to be confirmed after first clean run) \\
\midrule
Config hash   & \texttt{0000000000000000000000000000000000000000000000000000000000000001} \\
Build hash    & \texttt{0000000000000000000000000000000000000000000000000000000000000002} \\
Analysis hash & \texttt{0000000000000000000000000000000000000000000000000000000000000003} \\
Report hash   & \texttt{0000000000000000000000000000000000000000000000000000000000000004} \\
\bottomrule
\end{tabular}
\end{table}

\noindent These placeholder hashes will be replaced with confirmed values
from the first clean run of the full pipeline.
The \texttt{examples/vermont-2020-walkthrough/pin.sh} script pins the hashes
for the Vermont fixture after its first run.
